\section{Condiciones de la contratación}\label{sec:condiciones}
\subsection{Objeto, paquetes y valor referencial}
La contratación tiene por objeto el suministro, la instalación, la configuración, las pruebas y la puesta en servicio del servicio de información con terminales interactivos y alimentación solar híbrida de la EPIIS. Se organiza en cuatro paquetes para que cada proveedor acredite experiencia en su especialidad. Las especificaciones ET-01 a ET-10 son el requerimiento técnico y no mencionan marca \cite{compatibilizacion}.

\cuadro{tab:valor_referencial}{Valor referencial por paquete de contratación}{@{}P{1.2cm}Y P{3.3cm}r r@{}}{\textbf{Paq.}&\textbf{Alcance}&\textbf{Especificaciones}&\textbf{Base S/}&\textbf{Con contingencia S/}}{
P1 & Tres terminales interactivos con licencia, transporte a Cusco, anclaje y prueba & ET-01 & 35\,879,96 & 39\,467,96\\*
P2 & Generación FV, inversor, batería, tablero crítico, protecciones, conductores, tomacorrientes y puesta a tierra & ET-05 a ET-10 & 19\,306,21 & 21\,236,83\\
P3 & Cableado estructurado de siete enlaces, gabinete, switch y videovigilancia & ET-03 y ET-04 & 4\,419,00 & 4\,860,90\\
P4 & Ingeniería, configuración del CMS, interfaz, pruebas, documentación, formación y coordinación & ET-02 y plan de pruebas & 25\,320,00 & 27\,852,00\\\midrule
\textbf{Total} & & & \textbf{84\,925,17} & \textbf{93\,417,69}\\
}
Los importes provienen de los cuadros \ref{tab:precios_energia} a \ref{tab:presupuesto} e incluyen las reservas tributarias y de licencia de cada paquete. El valor referencial de la convocatoria se actualizará con la indagación formal de mercado y las cotizaciones que reciba la entidad.

\subsection{Plazo y lugar de ejecución}
El plazo total es de 112 días calendario, equivalentes a las 16 semanas del cronograma, contados desde el día siguiente de la notificación de la orden o la suscripción del contrato. Los bienes de los paquetes P1, P2 y P3 se entregarán en un máximo de 42 días calendario; su instalación se ejecutará entre las semanas 7 y 10 y las pruebas integradas entre las semanas 11 y 16.

La entrega e instalación se realizarán en el pabellón de la EPIIS, Ciudad Universitaria de Perayoc, Cusco: un terminal en cada piso, el gabinete y el tablero crítico en el tercer piso y los módulos en la cubierta. Transporte, descarga, ascenso y retiro de embalajes corren por cuenta del contratista.

\subsection{Forma de pago y conformidad}
\cuadro{tab:pagos}{Forma de pago por paquete}{@{}P{1.6cm}P{4.6cm}Y@{}}{\textbf{Paquete}&\textbf{Pago}&\textbf{Condición y documentos}}{
P1, P2 y P3 & Pago único por paquete & Conformidad de recepción, instalación y pruebas de su alcance, con acta de la puerta G3, guías de remisión, series, fichas, garantías y factura\\*
P4 & 30\% en G2; 40\% en G4; 30\% en G5 & Conformidad de los entregables E2, E4 y E5: diseño aprobado; servicio integrado con casos críticos superados; transferencia con manuales y formación\\
}
La conformidad la otorga la Dirección de la EPIIS como área usuaria, con informe técnico de TI para el software y la red y del responsable de infraestructura para la instalación eléctrica. El plazo para otorgarla y para pagar será el que fije el Reglamento de la Ley N.° 32069 \cite{contrataciones}. Las observaciones se notificarán por escrito con el plazo de subsanación y no se pagará un entregable observado.

\subsection{Requisitos del proveedor}
\cuadro{tab:requisitos_postor}{Requisitos de calificación del proveedor}{@{}P{3.0cm}Y@{}}{\textbf{Requisito}&\textbf{Condición y acreditación}}{
Capacidad legal & Inscripción vigente en el Registro Nacional de Proveedores para bienes o servicios, según el paquete; RUC activo y habido; declaración jurada de no estar impedido de contratar con el Estado\\*
Experiencia del postor & Facturación acumulada de una vez el valor referencial del paquete en bienes o servicios iguales o similares durante los ocho años anteriores, con contratos, órdenes o comprobantes de pago. Similares: pantallas o quioscos interactivos para P1; sistemas FV con baterías o instalaciones eléctricas para P2; cableado estructurado o videovigilancia para P3; implantación de sistemas o portales para P4\\
Personal clave de P2 & Ingeniero electricista o mecánico electricista colegiado y habilitado, con participación en al menos dos sistemas FV con baterías\\
Personal clave de P3 & Técnico en cableado estructurado con constancia o certificado de capacitación y experiencia en al menos dos instalaciones\\
Personal clave de P4 & Ingeniero de sistemas o informático con experiencia en al menos dos implantaciones de CMS, señalización digital o portales web\\
Seguridad del personal & Seguro Complementario de Trabajo de Riesgo para quienes trabajen en el pabellón y la cubierta; equipo de protección personal y procedimiento de trabajo en altura\\
Documentos técnicos & Fichas del fabricante que acrediten cada mínimo de las especificaciones; carta de garantía; declaración del canal de soporte en el Perú\\
}

\subsection{Garantía comercial y soporte}
\cuadro{tab:garantias}{Garantía comercial mínima y tiempos de atención}{@{}P{4.0cm}P{3.2cm}Y@{}}{\textbf{Bien o servicio}&\textbf{Garantía mínima}&\textbf{Atención}}{
Terminal interactivo & 1 año en sitio en Cusco & Respuesta en 48 h; solución o equipo de reemplazo en 10 días hábiles\\*
Inversor y batería & 5 años del fabricante & Diagnóstico remoto en 48 h; reemplazo en 15 días hábiles\\
Módulos FV & 12 años de producto y 25 de potencia & Trámite con el fabricante a cargo del proveedor\\
Cámaras, NVR, disco y switch & 1 año & Respuesta en 72 h; reemplazo en 10 días hábiles\\
Cableado estructurado & 1 año por instalación & Corrección de enlaces que no superen las pruebas, sin costo\\
Servicios de P4 & 90 días de soporte posterior a G5 & Incidencias críticas en 8 h de atención, conforme al RTO del servicio\\
}

\subsection{Penalidades}
La penalidad por mora se aplicará por cada día de atraso injustificado, con la fórmula y el tope que establece el Reglamento de la Ley N.° 32069 \cite{contrataciones}. Se establecen además otras penalidades objetivas y razonables, que se descontarán del pago o de la garantía de fiel cumplimiento:

\cuadro{tab:penalidades}{Otras penalidades aplicables}{@{}P{0.6cm}Y P{4.0cm}P{3.4cm}@{}}{\textbf{N.}&\textbf{Supuesto}&\textbf{Monto}&\textbf{Verificación}}{
1 & No atender una falla dentro del tiempo de respuesta de garantía & 0,5\% del valor del bien afectado por día & Registro de incidencias\\*
2 & Reemplazar personal clave sin aprobación previa & S/ 500 por ocurrencia & Informe del área usuaria\\
3 & No entregar planos de lo instalado, manuales o fichas en G5 & S/ 200 por día & Acta de transferencia\\
4 & Trabajar sin seguro complementario o sin equipo de protección & S/ 500 por ocurrencia & Informe del responsable de infraestructura\\
5 & Instalar un bien distinto del ofertado sin autorización & S/ 1\,000 por ocurrencia y retiro del bien & Acta de recepción\\
}
La suma de las otras penalidades no superará el tope que fije el Reglamento. Alcanzado el tope de cualquiera de ellas, la entidad podrá resolver el contrato.

\subsection{Responsabilidad por vicios ocultos}
Se consideran vicios ocultos los defectos que no eran apreciables en la recepción y que impiden el uso previsto de un bien o reducen su valor; por ejemplo, celdas de batería defectuosas, fallas intermitentes de la pantalla táctil o enlaces de red que se degradan. El contratista responde por ellos durante un año contado desde la conformidad, o durante el plazo de la garantía comercial si es mayor.

El procedimiento es el siguiente:
\begin{enumerate}
\item La entidad notifica el vicio por escrito dentro de los 30 días siguientes a su descubrimiento.
\item El contratista lo inspecciona dentro de los 5 días hábiles siguientes.
\item El contratista lo corrige o reemplaza el bien dentro de 15 días hábiles, sin costo para la universidad.
\end{enumerate}
Si no lo hace, la entidad podrá encargar la corrección a un tercero y repetir el costo contra el contratista, sin perjuicio de las acciones que correspondan.

\subsection{Responsabilidades del contratista}
\begin{itemize}
\item Suministrar bienes nuevos, del modelo ofertado y conformes con las especificaciones, con fichas, certificados y garantías.
\item Instalar, configurar y probar según el diseño aprobado, el Código Nacional de Electricidad y las buenas prácticas de cableado; ninguna modificación al diseño se ejecutará sin aprobación \cite{cne}.
\item Mantener el orden, la limpieza y la seguridad del área de trabajo, coordinar cortes de energía y no afectar servicios ajenos al piloto.
\item Entregar planos de lo instalado, unifilar, metrado final, registros de pruebas, manuales, inventario con series y credenciales bajo custodia.
\item Capacitar a editores, TI y operación local conforme al plan de formación y certificar la participación.
\item Guardar confidencialidad de credenciales, configuraciones y grabaciones; tratar los datos personales sólo para la finalidad del servicio \cite{datos}.
\item Retirar embalajes y residuos y entregar los aparatos o baterías reemplazados a un gestor autorizado \cite{raee}.
\end{itemize}

\subsection{Supervisión y control}
La Dirección de la EPIIS designará un coordinador del contrato, que verificará el avance contra el cronograma, aprobará los cambios menores y registrará las observaciones. TI supervisará la configuración, la seguridad y las pruebas del software y la red. El responsable de infraestructura supervisará la instalación eléctrica, la cubierta y los anclajes. Cada puerta de avance, de G1 a G5, se cierra con un acta firmada.

Los bienes recibidos se darán de alta en el inventario patrimonial de la universidad, con su ubicación, serie, garantía y responsable. Cualquier falla durante la ejecución se documentará y la resolverá el especialista que designe la autoridad competente.
